\section{Methodology}
\label{sec:methods}

\subsection{Empirical design and forecast metrics}

The analysis proceeds through four layers: point forecasting, forecast uncertainty, economic decision-making and tail-risk evaluation. Development choices are completed before the 2026 holdout is opened.

The naive benchmarks are
\begin{equation}
\widehat P^{(24)}_t=P_{t-24}
\end{equation}
and
\begin{equation}
\widehat P^{(168)}_t=P_{t-168}.
\end{equation}

Point forecasts are evaluated using MAE, RMSE and median absolute error. MAE is the primary model-selection criterion:
\begin{equation}
MAE=\frac{1}{N}\sum_{t=1}^{N}|P_t-\widehat P_t|.
\end{equation}
Percentage-error measures are not used because electricity prices can be zero or negative.

Every model within an outer fold is evaluated on an identical common-row set. Stress diagnostics additionally examine negative prices and realized prices above \EUR{200}/MWh and \EUR{500}/MWh.

\subsection{ElasticNet benchmark}

ElasticNet provides the regularized linear benchmark. Predictors are standardized within a pipeline, ensuring that scaling is refitted inside each inner training fold. The accepted benchmark uses four-fold chronological hourly \texttt{TimeSeriesSplit} cross-validation, MAE scoring, 15 logarithmically spaced \(\alpha\) values from \(10^{-3}\) to \(10^2\), and
\[
\rho\in\{0.1,0.5,0.9,1.0\}.
\]
This procedure was frozen before the later XGBoost design and was not retrospectively altered after the delivery-day inner-CV issue described below was recognized.

\subsection{XGBoost model and day-aligned tuning}

XGBoost uses the same Full and Tier-1 predictor sets as ElasticNet. The estimator uses the absolute-error objective, histogram tree construction and random seed 42. The pre-specified grid is
\[
\text{max depth}\in\{3,4,5\},
\]
\[
\text{n estimators}\in\{200,400\},
\]
\[
\text{learning rate}\in\{0.03,0.05,0.10\},
\]
and
\[
\text{subsample}\in\{0.8,1.0\},
\]
giving 36 configurations.

XGBoost hyperparameters are selected using three chronological inner folds over unique Europe/Berlin \emph{delivery dates}. This keeps all 23, 24 or 25 observations of a local delivery day on the same side of an inner train-validation boundary. The choice matches the economic forecast origin: all prices for day \(D\) are formed together at D-1 11:45, so realized targets from the early hours of \(D\) should not influence model selection for later hours of that same day.

The development promotion rule required Full XGBoost to improve on the frozen Full ElasticNet pipeline by more than 1\% in row-weighted outer-fold MAE. The rule was specified before the first real-data XGBoost result.

\subsection{Final refit}

Once the development procedure was complete, Full and Tier-1 XGBoost were refitted independently on all development data. Each information set reran the same frozen 36-configuration search; hyperparameters were not transferred from one information set to the other.

The selected final configurations were:

\begin{table}[htbp]
\centering
\caption{Frozen final XGBoost configurations.}
\label{tab:final_hyperparams}
\begin{adjustbox}{max width=\textwidth}
\begin{tabular}{lrrrrr}
\toprule
Information set & Learning rate & Depth & Trees & Subsample & Rows \\
\midrule
Full & 0.10 & 3 & 200 & 0.80 & 61,142 \\
Tier-1 & 0.03 & 3 & 400 & 0.80 & 61,143 \\
\bottomrule
\end{tabular}
\end{adjustbox}
\end{table}

Model files, predictor lists and manifests were checksum-verified before holdout evaluation.

\subsection{Delivery-day-safe empirical uncertainty}

Uncertainty is constructed from genuinely out-of-sample development residuals,
\begin{equation}
e_t=P_t-\widehat P_t.
\end{equation}
The outer-fold XGBoost predictions from 2023--2025 form a continuous walk-forward residual history.

For delivery day \(D\), only residuals from delivery dates strictly before \(D\) are eligible. The empirical residual quantile is
\begin{equation}
r_D(q)=Q_q\left(\{e_t: D-w\leq delivery\_date(t)<D\}\right),
\end{equation}
and the corresponding hourly price quantile is
\begin{equation}
\widehat Q_{D,h}(q)=\widehat P_{D,h}+r_D(q).
\end{equation}

The frozen quantiles are \(q\in\{0.10,0.50,0.90\}\). Thus the nominal central 80\% interval is
\begin{equation}
[L_{D,h},U_{D,h}]
=
[\widehat Q_{D,h}(0.10),\widehat Q_{D,h}(0.90)].
\end{equation}

Candidate residual windows are \(w\in\{60,90,120,180,365\}\) delivery days. Each requires at least
\[
\max(1,\lfloor w/4\rfloor)
\]
days of history. Candidates are evaluated on one common observation set using the Winkler interval score. With \(\alpha=0.20\),
\begin{align}
IS_\alpha(L,U;P)
&=(U-L)
+\frac{2}{\alpha}(L-P)\mathbf{1}(P<L) \nonumber\\
&\quad+\frac{2}{\alpha}(P-U)\mathbf{1}(P>U).
\end{align}
Lower is better. The window with the lowest common-row score is selected mechanically; coverage diagnostics cannot override that result.

\subsection{Economic decision rule}

The economic layer is a stylized single-cycle daily storage problem. For each delivery day, all pairs \((i,j)\) satisfying \(i<j\) are candidates. A normalized one-unit charge is made at \(i\), with round-trip efficiency \(\eta_{rt}\) applied to the discharge leg.

Throughput degradation cost is
\begin{equation}
C_{deg}=c(1+\eta_{rt}),
\end{equation}
and market transaction cost is frozen at
\[
C_{market}=0.
\]
Total cost is
\begin{equation}
C=C_{deg}+C_{market}.
\end{equation}

For a point forecast, the decision score is
\begin{equation}
Score_P(i,j)=\eta_{rt}\widehat P_j-\widehat P_i-C.
\end{equation}
The strategy selects the highest-scoring pair and abstains if the best score is non-positive.

After the pair is fixed, realized net P\&L is
\begin{equation}
\Pi_D=\eta_{rt}P_{D,j^*}-P_{D,i^*}-C.
\end{equation}
Realized prices therefore value a decision already made from forecasts; they never select the trade.

For uncertainty-aware strategies, the conservative score is
\begin{equation}
Score_U(i,j)=\eta_{rt}L_j-U_i-C.
\end{equation}
This is a decision filter based on marginal bounds, not a jointly calibrated interval for the price spread.

The frozen strategies are:
\begin{description}[leftmargin=1.2cm,style=nextline]
\item[S0] no trade;
\item[S1] lag-24 point forecast;
\item[S2] XGBoost Full point forecast;
\item[S3] XGBoost Full with uncertainty filter;
\item[S4] XGBoost Tier-1 point forecast;
\item[S5] XGBoost Tier-1 with uncertainty filter.
\end{description}

The four pre-specified contrasts are
\begin{align}
\Delta_{forecast}&=\Pi_{S2}-\Pi_{S1},\\
\Delta_{uncertainty}&=\Pi_{S3}-\Pi_{S2},\\
\Delta_{Tier1}&=\Pi_{S4}-\Pi_{S2},\\
\Delta_{Tier1,U}&=\Pi_{S5}-\Pi_{S3}.
\end{align}

Economic comparisons use only delivery days on which every strategy has a complete intraday vector.

\subsection{Sensitivity, oracle and tail risk}

The primary parameters are
\[
\eta_{rt}=0.85,\qquad c=\EUR{10}.
\]
A complete sensitivity grid is frozen in advance:
\[
\eta_{rt}\in\{0.70,0.85,0.92\},\qquad
c\in\{5,10,15\}.
\]

The primary holdout criterion concerns only \(\Delta_{forecast}\). Full confirmation requires a positive value in the primary cell and all nine sensitivity cells; partial confirmation requires a positive primary cell but fewer than nine positive cells; failure occurs if the primary cell is non-positive.

An ex-post perfect-information oracle is
\begin{equation}
\Pi^{oracle}_D=
\max\left[0,\max_{i<j}(\eta_{rt}P_j-P_i-C)\right].
\end{equation}
It is isolated from the executable strategy and used only as an upper-bound diagnostic. Value capture is
\begin{equation}
VCR_s=\frac{\sum_D\Pi_{D,s}}{\sum_D\Pi^{oracle}_D}.
\end{equation}

Daily strategy loss is defined as
\begin{equation}
L_{D,s}=-\Pi_{D,s}.
\end{equation}
Empirical historical VaR and Expected Shortfall are reported at 95\% and 99\%. ES is computed from the exact worst \(m=(1-c)N\) observation-equivalents, using fractional weight at the boundary when necessary. Estimates based on fewer than 20 effective tail observations are flagged as low precision.

\subsection{Holdout updating and frozen protocol}

At the start of 2026, uncertainty history consists of the development out-of-sample residuals, not in-sample residuals from the final refitted models. For each holdout day \(D\), the uncertainty snapshot is formed using eligible development residuals plus residuals from already revealed holdout days. The day-\(D\) residual is appended only after prices for \(D\) are revealed and can therefore affect \(D+1\) onward.

The holdout boundary, final model-refit rule, uncertainty specification, strategy definitions, primary parameter cell, nine-cell grid, common-day rule, VaR/ES method and mechanical confirmation rule were all frozen before the first 2026 outcome metric was inspected. The complete holdout orchestration was rehearsed on non-holdout data before exposure. If a bug had been discovered after exposure, the protocol required preservation of the original result and explicit labeling of any corrective analysis rather than silent replacement.
